\subsection{Wortinduktion in Konkatenationsschreibweise}

Die Minimalität des konkreten Wortmodells wurde vor Einführung der Axiome
aus seiner Konstruktion als kleinste abgeschlossene Menge bewiesen.
Jetzt wenden wir das
axiomatische Induktionsprinzip W3 auf das konkrete Wortmodell an
und schreiben den Anfügungsschritt als Konkatenation.

\noindent\begin{minipage}{\linewidth}
\FormulaThmDeltaK[Strukturinduktion über alle endlichen Wörter]{%
  \begin{aligned}[t]
    &P(\EmptyWord A)\dsep{}\\[-2pt]
    &\forall u\in\FiniteWords A\,\forall a\in A\,
      (P(u)\rightarrow P(\WordConcat u{\LetterWord a}))\\[-2pt]
    &\vdash\forall w\in\FiniteWords A\,P(w)
  \end{aligned}
}{FiniteWordStructuralInduction}{%
  \DeltaRow{Mengen}{A\dsep u\dsep w\dsep a\dsep n}
  \DeltaRow{Prädikat}{P}
}
\end{minipage}\par
Wir wenden \FormulaRefAuto{WordStructureInduction}[ax] auf das Modell
\FormulaRefAuto{WordStructureConcreteModel}[thm] an. Die Brücke
\FormulaRefAuto{WordModelSuccessorConcatenation}[thm] übersetzt den
primitiven Nachfolger in das Anhängen eines Buchstabenwortes.
\begin{tabproofwide}
  \proofstepwidestar[1]{P(\EmptyWord A)}{\rA}
  \proofstepwidestar[2]{\forall u\in\FiniteWords A\,\forall a\in A\;(P(u)\rightarrow P(\WordConcat u{\LetterWord a}))}{\rA}
  \proofstepwidestar[]{\mathsf{WortStr}_A(\FiniteWords A,\EmptyWord A,\sigma_A)}{\FormulaRefAuto{WordStructureConcreteModel}[thm]}
  \proofstepwidestar[4]{u\in\FiniteWords A}{\rA}
  \proofstepwidestar[5]{a\in A}{\rA}
  \proofstepwidestar[2,4,5]{P(u)\rightarrow P(\WordConcat u{\LetterWord a})}{\rRE{5,\rUE{\rRE{4,\rUE{2}}}}}
  \proofstepwidestar[4,5]{\sigma_A(u,a)=\WordConcat u{\LetterWord a}}{\FormulaRefAuto{WordModelSuccessorConcatenation}[thm]{4,5}}
  \proofstepwidestar[2,4,5]{P(u)\rightarrow P(\sigma_A(u,a))}{\rIE{7,6}}
  \proofstepwidestar[2]{\forall u\in\FiniteWords A\,\forall a\in A\;(P(u)\rightarrow P(\sigma_A(u,a)))}{\rUI{\rRI{4,\rUI{\rRI{5,8}}}}}
  \proofstepwidestar[1,2]{\forall w\in\FiniteWords A\;P(w)}{\FormulaRefAuto{WordStructureInduction}[thm]{3,1,9}}
\end{tabproofwide}


\noindent\begin{minipage}{\linewidth}
\FormulaThmDeltaK[Die Wortmenge ist die kleinste unter Anfügung abgeschlossene Menge]{%
  \begin{aligned}[t]
    &\EmptyWord A\in C\dsep{}\\[-2pt]
    &\forall u\in C\cap\FiniteWords A\,\forall a\in A\,
       (u\cup\{(\WordLength u,a)\}\in C)\\[-2pt]
    &\vdash\FiniteWords A\subseteq C
  \end{aligned}
}{FiniteWordSetLeastAppendClosed}{%
  \DeltaRow{Mengen}{A\dsep C\dsep u\dsep a\dsep w}
}
\end{minipage}\par
Die Schnittmenge in der Voraussetzung stellt lediglich sicher, dass die
Wortlänge im Anfügungsterm definiert ist. Man kann insbesondere
\(C\subseteq\FiniteWords A\) voraussetzen; dann folgt \(C=\FiniteWords A\).
Zusammen mit dem Leerwort und der abgeschlossenen Anfügung ist dies die
gewünschte Erzeugungsregel für \(A^*\).
\begin{tabproofwide}
  \proofstepwidestar[1]{\EmptyWord A\in C}{\rA}
  \proofstepwidestar[2]{\forall u\in C\cap\FiniteWords A\,\forall a\in A\,
    (u\cup\{(\WordLength u,a)\}\in C)}{\rA}
  \proofstepwidestar[3]{u\in\FiniteWords A}{\rA}
  \proofstepwidestar[4]{a\in A}{\rA}
  \proofstepwidestar[5]{u\in C}{\rA}
  \proofstepwidestar[3,5]{u\in C\cap\FiniteWords A}{%
    \FormulaRefAuto{x\in A\cap B\eqvdash x\in A\land x\in B}[thm]{\rAI{5,3}}}
  \proofstepwidestar[2,3,4,5]{u\cup\{(\WordLength u,a)\}\in C}{\rRE{4,\rUE{\rRE{6,\rUE{2}}}}}
  \proofstepwidestar[3,4]{\WordConcat u{\LetterWord a}=u\cup\{(\WordLength u,a)\}}{%
    \FormulaRefAuto{WordAppendOccurrenceEquation}[thm]{3,4}}
  \proofstepwidestar[2,3,4,5]{\WordConcat u{\LetterWord a}\in C}{\rIE{\FormulaRefAuto{a=b\vdash b=a}[thm]{8},7}}
  \proofstepwidestar[2]{\forall u\in\FiniteWords A\,\forall a\in A\,
    (u\in C\rightarrow\WordConcat u{\LetterWord a}\in C)}{\rUI{\rRI{3,\rUI{\rRI{4,\rRI{5,9}}}}}}
  \proofstepwidestar[1,2]{\forall w\in\FiniteWords A\,(w\in C)}{\FormulaRefAuto{FiniteWordStructuralInduction}[thm]{1,10}}
  \proofstepwidestar[1,2]{\FiniteWords A\subseteq C}{%
    \FormulaRefAuto{A\subseteq B\coloneq\forall x\,(x\in A\rightarrow x\in B)}[def]{11}}
\end{tabproofwide}


